\section{Introduction to Heating \& Cooling Equipment}

\subsection{Overview}
Heating and cooling equipment are fundamental to industrial processes, commercial buildings, and residential applications. They involve the practical application of heat transfer and thermodynamic principles to either generate heat for various uses (e.g., steam production, space heating) or remove heat to achieve cooling (e.g., refrigeration, air conditioning).

\subsection{Basic Principles}
At its core, heating equipment (like boilers) converts chemical energy from fuels into thermal energy, which is then transferred to a working fluid (typically water to produce steam). Cooling equipment (like refrigerators, air conditioners, and chillers) operates on thermodynamic cycles to transfer heat from a lower-temperature region to a higher-temperature region, requiring work input.

\section{Heating Equipment: Boilers}

\subsection{What is a Boiler?}
A boiler (or steam generator) is a closed vessel in which water or other fluid is heated. The heated or vaporized fluid exits the boiler for use in various processes or heating applications. Boilers are designed to efficiently transfer heat from the combustion of fuels to the working fluid, producing hot water or steam at desired pressures and temperatures.

\subsection{Components and Mechanism}
Boilers consist of a furnace (combustion chamber) where fuel is burned, heat transfer surfaces (tubes or shell), and a steam drum (for separating steam from water). Heat is transferred from the hot combustion gases to the water primarily by radiation, convection, and conduction.

\subsection{Types of Boilers}
Boilers are classified based on their construction, fuel type, heat recovery method, and application:
\begin{enumerate}
    \item \textbf{Water Tube Boiler:}
        \begin{itemize}
            \item \textit{Mechanism:} Water circulates inside tubes, and hot combustion gases flow around the outside of the tubes.
            \item \textit{Characteristics:} Designed for high-pressure and high-temperature steam production, making them suitable for power generation plants. They offer better heat transfer rates and are safer at high pressures due to the smaller volume of water in the tubes.
        \end{itemize}
    \item \textbf{Fire Tube Boiler:}
        \begin{itemize}
            \item \textit{Mechanism:} Hot combustion gases pass through tubes, and water surrounds the outside of the tubes within a large shell.
            \item \textit{Characteristics:} Simpler in construction, generally used for lower-pressure and lower-temperature steam applications. Common in small-scale industrial applications and historic locomotives.
        \end{itemize}
    \item \textbf{Package Boiler:}
        \begin{itemize}
            \item \textit{Mechanism:} A factory-made, complete boiler system that is pre-assembled and shipped as a unit. Can be either fire-tube or water-tube.
            \item \textit{Characteristics:} Compact, easy to install, and designed for relatively low-pressure applications.
        \end{itemize}
    \item \textbf{Stocker Fired Boiler:}
        \begin{itemize}
            \item \textit{Mechanism:} Utilizes a mechanical stoker system to feed solid fuels (like coal, biomass, or municipal solid waste) onto a grate in the furnace. Air is supplied from below for combustion.
            \item \textit{Characteristics:} Efficient for burning a wide range of solid fuels.
        \end{itemize}
    \item \textbf{Pulverized Coal Boiler:}
        \begin{itemize}
            \item \textit{Mechanism:} Coal is ground into a fine powder (pulverized) and then blown into the furnace, where it burns in suspension like a gas.
            \item \textit{Characteristics:} Allows for complete and efficient combustion of coal, leading to high thermal efficiencies. Widely used in large thermal power plants.
        \end{itemize}
    \item \textbf{Waste Heat Boiler (WHB):}
        \begin{itemize}
            \item \textit{Mechanism:} Recovers heat from hot flue gases or exhaust air streams from industrial processes (e.g., gas turbines, furnaces) to generate steam or hot water.
            \item \textit{Characteristics:} Improves overall energy efficiency by utilizing otherwise wasted energy, reducing fuel costs and emissions.
        \end{itemize}
    \item \textbf{Fluidized Bed Boiler (FBB) / Circulating Fluidized Bed (CFB) Boiler:}
        \begin{itemize}
            \item \textit{Mechanism:} Fuel particles are suspended in a bed of inert material (like sand or ash) by an upward flow of air. Combustion occurs within this turbulent mixture, allowing for efficient mixing and heat transfer.
            \item \textit{Characteristics:} Offers very high combustion efficiency, can burn a variety of fuels (including low-grade coals and biomass), and is excellent for controlling emissions (e.g., SOx and NOx) due to lower combustion temperatures and effective sorbent injection. CFB boilers involve a circulating loop of solids for even better performance.
        \end{itemize}
\end{enumerate}

\section{Cooling Equipment: Refrigeration Cycle Components \& Systems}

\subsection{The Cooling Process Overview}
The fundamental cooling process in refrigeration involves a closed-loop system where a refrigerant circulates through four main components: a compressor, a condenser, an expander (or throttling device), and an evaporator. This cycle continually absorbs heat from a low-temperature space and rejects it to a higher-temperature environment.

\subsection{Core Components of Refrigeration Cycle}
\begin{enumerate}
    \item \textbf{Compressor:}
        \begin{itemize}
            \item \textit{Function:} The "heart" of the refrigeration cycle. It draws in low-pressure, low-temperature refrigerant vapor from the evaporator and compresses it to a high-pressure, high-temperature vapor. This increases the saturation temperature of the refrigerant, allowing it to condense at a higher temperature than the surroundings.
            \item \textit{Types:} Common types include:
                \begin{itemize}
                    \item \textit{Reciprocating/Piston Compressors:} Use pistons to compress the refrigerant.
                    \item \textit{Rotary-Screw Compressors:} Use two helical rotors to compress the gas.
                    \item \textit{Scroll Compressors:} Use two interleaved scrolls, one fixed and one orbiting, to compress the gas.
                \end{itemize}
        \end{itemize}
    \item \textbf{Condenser:}
        \begin{itemize}
            \item \textit{Function:} Rejects heat from the high-pressure, high-temperature refrigerant vapor to a cooling medium (typically ambient air or water). As heat is removed, the refrigerant condenses into a high-pressure liquid.
            \item \textit{Types:} Classified by the cooling medium:
                \begin{itemize}
                    \item \textit{Air-Cooled Condensers:} Use fans to blow air over the condenser coils. Common in residential air conditioners.
                    \item \textit{Water-Cooled Condensers:} Use cooling water (often from a cooling tower) to absorb heat. Common in large commercial and industrial systems.
                    \item \textit{Evaporative Condensers:} Combine aspects of air-cooled and water-cooled, using both air and evaporative cooling of water to reject heat.
                \end{itemize}
        \end{itemize}
    \item \textbf{Expander (Expansion Valve / Throttling Device):}
        \begin{itemize}
            \item \textit{Function:} Reduces the pressure of the high-pressure liquid refrigerant, causing a drop in its temperature. This prepares the refrigerant for heat absorption in the evaporator. The expansion process is typically isenthalpic (constant enthalpy).
            \item \textit{Types:}
                \begin{itemize}
                    \item \textit{Thermostatic Expansion Valve (TXV/TSV):} Regulates the flow of refrigerant to the evaporator based on the superheat of the refrigerant vapor leaving the evaporator.
                    \item \textit{Capillary Tube:} A long, thin tube that provides a fixed restriction to refrigerant flow, relying on pressure drop along its length to achieve expansion. Simple and inexpensive, often used in household refrigerators.
                \end{itemize}
        \end{itemize}
    \item \textbf{Evaporator:}
        \begin{itemize}
            \item \textit{Function:} Where the low-pressure, low-temperature liquid-vapor mixture of refrigerant absorbs heat from the refrigerated space or fluid to be cooled, causing the refrigerant to evaporate completely into a low-pressure vapor.
            \item \textit{Refrigerants:} Common refrigerants used include various hydrofluorocarbons (HFCs) like R-134a, or natural refrigerants like ammonia and CO$_2$.
        \end{itemize}
\end{enumerate}

\subsection{Types of Cooling Equipment/Systems}
\begin{itemize}
    \item \textbf{Air Conditioners:} Remove heat and humidity from indoor air and reject it outdoors, providing space cooling.
    \item \textbf{Chillers:} Use a refrigeration cycle to cool a secondary fluid (typically water or a brine solution), which is then circulated to absorb heat from specific areas or processes. Widely used in commercial buildings and industrial facilities.
    \item \textbf{Cooling Towers:} Reject waste heat to the atmosphere, primarily through the evaporation of a small portion of circulating water. They are used in conjunction with water-cooled condensers in large refrigeration and power generation systems.
    \item \textbf{Cold Rooms / Refrigerated Warehouses:} Specialized, insulated rooms designed to maintain low temperatures for the storage of perishable goods or other temperature-sensitive items.
    \item \textbf{Refrigerators (Household):} Compact appliances designed for domestic food preservation.
\end{itemize}

\section{Performance Measures}
The performance of heating and cooling equipment is primarily evaluated by their efficiency in converting energy input into useful heat or cooling output.
\begin{itemize}
    \item \textbf{Boiler Efficiency:} For boilers, this is typically the ratio of the useful heat transferred to the working fluid (steam or hot water) to the chemical energy input from the fuel. (Specific formulas and calculations are covered in the Thermodynamics or Heat Exchanger chapters).
    \item \textbf{Coefficient of Performance (COP):} For refrigeration and heat pump systems, performance is measured by COP (as detailed in the "Refrigeration" chapter), which quantifies the ratio of desired heat transfer to work input.
\end{itemize}

\section{Formulas Summary for Equipment for Heating \& Cooling}

\subsection{Boiler Heat Output}
The heat output of a boiler (rate of heat transferred to the working fluid) can be expressed as:
\begin{equation}
    \Qdot_{boiler} = \dot{m}_{steam} (h_{steam,out} - h_{water,in})
\end{equation}
where: \\
$\Qdot_{boiler}$ = Rate of heat output from the boiler (W) \\
$\dot{m}_{steam}$ = Mass flow rate of steam/hot water produced (kg/s) \\
$h_{steam,out}$ = Enthalpy of steam/hot water leaving the boiler (J/kg) \\
$h_{water,in}$ = Enthalpy of feed water entering the boiler (J/kg)

\subsection{Heat Input to Boiler (from Fuel)}
\begin{equation}
    \Qdot_{fuel,in} = \dot{m}_{fuel} \times \text{Heating Value of Fuel}
\end{equation}
where: \\
$\Qdot_{fuel,in}$ = Rate of energy input from fuel (W) \\
$\dot{m}_{fuel}$ = Mass flow rate of fuel (kg/s) \\
Heating Value of Fuel = Energy released per unit mass of fuel (J/kg) (e.g., HHV - Higher Heating Value)

\subsection{Boiler Thermal Efficiency}
\begin{equation}
    \eta_{boiler} = \frac{\Qdot_{boiler}}{\Qdot_{fuel,in}} \times 100\%
\end{equation}
where: \\
$\eta_{boiler}$ = Boiler thermal efficiency (dimensionless or percentage)

\subsection{Refrigeration Cycle Component Energy Balances}
(These fundamental energy balance equations for compressors, condensers, evaporators, and expansion valves are detailed in the "Refrigeration" chapter.)
